\begin{quote}
\small
\noindent\textbf{Resumo.} A Geração Aumentada por Recuperação (RAG) é a principal forma de ancorar respostas de modelos de linguagem em documentos, mas a escolha do método de injeção de contexto afeta custo e qualidade de maneiras pouco medidas fora do inglês. Este trabalho replica, com documentos em português, a comparação de seis métodos de RAG (Stuff, Refine, Map-Reduce, Map-Rerank, Query Step-Down e Reciprocal RAG) proposta por Şakar e Emekci (2025), em três domínios brasileiros com respostas de referência verificáveis: a prova discursiva do Revalida 2024/2, as Informações Trimestrais da Grendene S.A. e um artigo científico sobre RAG em português. As 162 respostas foram avaliadas por similaridade de cosseno e por quatro métricas RAGAS, com validação do juiz automático contra ele mesmo e contra um juiz de outra família de modelos. Também foram avaliados filtros de compressão contextual e o custo de infraestrutura de ChromaDB e FAISS. Nenhum método foi estatisticamente superior em similaridade de cosseno, enquanto o custo variou até cinco vezes entre eles. A ordenação dos métodos mudou com a métrica: o Map-Rerank teve a maior similaridade de cosseno e uma das menores \emph{Answer Relevancy}. O método mais simples, o Stuff, recusou 30\% das perguntas, na maioria das vezes com a informação presente no contexto recuperado. Limiares fixos de compressão ou não tiveram efeito, ou descartaram quase todo o contexto, enquanto limiares calibrados na distribuição de similaridade reduziram cerca de 40\% dos tokens sem perda de qualidade em dois dos três domínios. Os resultados indicam que a escolha do método em RAG é sobretudo uma decisão de custo e do tipo de resposta esperada, e que conclusões baseadas em uma única métrica automática devem ser vistas com cautela.

\medskip
\noindent\textbf{Palavras-chave:} geração aumentada por recuperação; modelos de linguagem; avaliação de RAG; RAGAS; compressão contextual; português.

\bigskip
\noindent\textbf{Abstract.} Retrieval-Augmented Generation (RAG) is the main way to ground language model answers in documents, but the choice of context injection method affects cost and quality in ways rarely measured outside English. This work replicates, with Portuguese documents, the comparison of six RAG methods (Stuff, Refine, Map-Reduce, Map-Rerank, Query Step-Down and Reciprocal RAG) proposed by Şakar and Emekci (2025), in three Brazilian domains with verifiable reference answers: the Revalida 2024/2 medical licensing exam, the quarterly financial statements of Grendene S.A., and a scientific paper on RAG in Portuguese. The 162 answers were evaluated with cosine similarity and four RAGAS metrics, and the automatic judge was validated against itself and against a judge from another model family. Contextual compression filters and the infrastructure cost of ChromaDB and FAISS were also evaluated. No method was statistically superior in cosine similarity, while cost varied up to fivefold across methods. The ranking of methods changed with the metric: Map-Rerank had the highest cosine similarity and one of the lowest Answer Relevancy scores. The simplest method, Stuff, refused to answer 30\% of the questions, mostly when the information was present in the retrieved context. Fixed compression thresholds either had no effect or discarded almost all context, whereas thresholds calibrated on the similarity distribution cut about 40\% of tokens without quality loss in two of the three domains. The results indicate that choosing a RAG method is mainly a decision about cost and the expected answer type, and that conclusions based on a single automatic metric should be taken with caution.

\medskip
\noindent\textbf{Keywords:} retrieval-augmented generation; language models; RAG evaluation; RAGAS; contextual compression; Portuguese.
\end{quote}
